Подтверждающий эксперимент даёт положительный результат с определёнными границами: модель, обученная до тяжёлого кризиса, может снизить стоимость реакции относительно статического планирования и фиксированного правила добавления мощности в условии с обратной связью через возвраты. Выбор компонент по точности прогноза часто исключает важную зависимость возвратов. При исходной стоимости выбор действий для разных механизмов также улучшает описательную нормированную стоимость относительно обеих фиксированных стратегий. Однако методика не превосходит их во всех условиях и не гарантирует предотвращения кризиса. Проигрыш без обратной связи, чувствительность к стоимости действия и зависимость от заданных эффектов вмешательств остаются существенными ограничениями.
